\begin{introduction}\label{chapter:introduction}
% R02 en ../write_thesis_roadmap.md. Sin referencias nuevas: todas las claves están en
% Bibliography.bib desde R01. Hipótesis H1-H3 de docs/resumen.md; la forma de H2 y H3 la
% decidió Josue el 2026-11-01 (H2 se mantiene sin experimento; H3 como lectura de WSE-bench).

Generar historias de forma automática es uno de los problemas más antiguos de la inteligencia
artificial. TALE-SPIN~\cite{meehan1977}, a finales de los años setenta, ya producía fábulas
razonando cómo cada personaje resolvía sus objetivos. Durante décadas, los generadores dependieron
de mundos descritos a mano en lenguajes formales, que garantizaban la causalidad de los eventos a
costa de acotar las historias posibles~\cite{riedl2010}. Los \textit{Large Language Models}
(LLMs), modelos de lenguaje de gran tamaño capaces de seguir instrucciones en lenguaje natural,
cambiaron ese límite: producen prosa fluida a partir de una petición, sin un modelo del dominio
programado de antemano. Lo que no garantizan es que esa prosa sea una buena historia, y tampoco
ofrecen una forma fiable de saber si lo es. Esta tesis trata de esos dos problemas: generar
historias con LLMs y evaluarlas.

La literatura ordena los generadores según quién decide la trama. En la generación dirigida por la
trama (\textit{Top-Down}), un componente con visión de autor decide los eventos y los personajes
los ejecutan; con LLMs, eso se traduce en planificar la historia antes de
escribirla~\cite{yang2023,huot2025}. En la generación por simulación de personajes
(\textit{Bottom-Up}), la trama no existe de antemano: emerge de lo que hacen unos agentes con
objetivos y memoria propios~\cite{aylett1999,park2023}. Entre ambas hay una tradición de sistemas
híbridos, como Façade~\cite{mateas2005} o The Virtual Storyteller~\cite{theune2003}, que combinan
un componente de autor con personajes autónomos y que los LLMs han reabierto~\cite{chen2026a}. El
Capítulo~\ref{chapter:state-of-the-art} precisa estos términos y recorre las tres familias.

El antecedente directo de este trabajo es la tesis de Fuentes Rodríguez~\cite{fuentes2025},
presentada en esta facultad en 2025 y cotutorada por el mismo tutor. Su sistema generaba
historias mediante la simulación de agentes conversacionales y lo comparaba con un modelo guiado
por un solo \textit{prompt} (la instrucción en lenguaje natural que recibe el modelo). Doce
voluntarios puntuaron ambas salidas de 1 a 10, y los perfiles resultantes apenas se distinguieron.
De ese trabajo se conservan la línea base de un solo \textit{prompt} y el interés por la coherencia
de las historias largas; se cambia la forma de evaluar.

Los problemas que motivan esta tesis se agrupan en tres:

\begin{itemize}
  \item \textbf{La paradoja narrativa:} la simulación de personajes da comportamientos creíbles,
  pero nadie vela por la estructura que convierte una serie de eventos en una historia; la
  planificación vela por esa estructura, pero tiende a reducir a los personajes a ejecutores del
  plan~\cite{louchart2003,riedl2010}. Un generador que use solo una de las dos renuncia a lo que
  aporta la otra.
  \item \textbf{La pérdida de compromisos y de riqueza:} aun con el contexto necesario delante, los
  LLMs dejan sin pagar muchas de las prefiguraciones que introducen~\cite{yun2026}, se contradicen a
  lo largo de una historia larga~\cite{li2026} y hacen que un personaje revele lo que en ese momento
  no podía saber~\cite{ahn2024}. Razonar quién sabe qué en una conversación les resulta mucho más
  difícil que a las personas~\cite{kim2023}, y sus historias superan menos pruebas de creatividad
  que las de escritores profesionales~\cite{chakrabarty2024}.
  \item \textbf{La evaluación poco fiable:} las escalas de valoración absolutas son las más usadas y
  las más criticadas~\cite{amidei2019}, y con ellas incluso evaluadores expertos coinciden
  poco~\cite{chiang2023}. Los LLMs usados como jueces favorecen sus propias generaciones y las
  historias escritas por IA~\cite{zheng2023,panickssery2024,tutone2026}. Sin una evaluación fiable
  no se puede saber si un mecanismo nuevo mejora la historia o solo la cambia.
\end{itemize}

El tema es actual. Los sistemas híbridos con LLMs son de los últimos años: StoryBox~\cite{chen2026a}
y WSE-bench~\cite{chen2026b}, este último todavía un preprint, combinan agentes simulados con un
narrador, y Yu et al.~\cite{yu2025}, con un plan. También son recientes el trabajo de Yun et al.~\cite{yun2026} sobre los
compromisos sin pagar y los resultados que cuestionan a los jueces LLM en escritura creativa o que
proponen aprenderlos de preferencias humanas~\cite{tutone2026,fein2026}. Además, las historias de
esta tesis están en español, una lengua en la que, en una comparación de GPT-4 con un novelista, la
brecha entre ambos fue mayor que en inglés~\cite{marco2024}.

La \textbf{novedad científica} del trabajo está en dos aportaciones. La primera es un generador
híbrido, bidireccional en el sentido del título: el plan baja de arriba abajo
(\textit{Top-Down}) y los personajes simulados suben de abajo arriba (\textit{Bottom-Up}). Separa tres
decisiones. Un plan validado decide qué pasa; unos personajes simulados que
solo saben lo que presenciaron deciden cómo pasa, al representar el plan escena a escena en lo que
llamamos una \textbf{función}, como en el teatro; y un narrador que cura el registro de esa función,
sin inventar, decide cómo se cuenta. En la literatura que revisa el
Capítulo~\ref{chapter:state-of-the-art}, que no es exhaustiva, no encontramos un generador con LLMs
que reúna las tres piezas tal como aquí se formulan; los dos más cercanos, Yu et al. y PLOTTER,
comparten solo parte de ellas~\cite{yu2025,gu2026} (Sección~\ref{sec:posicion}). La segunda es una
metodología de evaluación en la que el LLM no juzga la calidad: cuenta rasgos de la historia y cita el
pasaje de cada recuento, que el código verifica literalmente. La calidad la deciden lectores que
comparan pares de historias a ciegas, y un juez aprendido de esas comparaciones intenta ordenar las
historias como ellos a partir de los recuentos, con una validación fuera de muestra que dice si lo
consigue.

La importancia teórica del trabajo es doble. Como el plan, la función y la narración son etapas
separadas, el generador sirve de banco de pruebas de la paradoja narrativa: se puede cambiar una
etapa y conservar las demás, por ejemplo escribir la historia directamente desde el mismo plan o
narrar la misma función desde otro punto de vista, y medir qué aporta cada una. La metodología, por
su parte, muestra cómo evaluar historias sin pedir a un modelo un juicio de calidad, una cuestión
abierta en el área (Sección~\ref{sec:evaluacion-historias}). En la práctica, el generador y las
herramientas de evaluación forman un sistema que recalcula sus resultados desde los datos guardados,
que residen en la copia del autor y no en Git (Sección~\ref{sec:pruebas}), con un bot de Telegram para pedir
historias y una interfaz local para leerlas a ciegas, y el juez lee solo el texto de la historia,
de modo que la metodología se puede aplicar a las de otro generador.

A partir de lo anterior se plantea el siguiente \textbf{problema}: ¿cómo generar historias con LLMs
que conserven los compromisos de su trama sin reducir a los personajes a ejecutores del plan, y cómo
evaluarlas de forma fiable sin pedir a un modelo que juzgue su calidad? El \textbf{objeto de
estudio} es la generación automática de historias con modelos de lenguaje, y el \textbf{campo de
acción}, la generación híbrida de historias en español, que combina la planificación de la trama
con la simulación de personajes, junto con su evaluación mediante rasgos contables y comparaciones
humanas por pares.

Para resolver el problema se plantea como \textbf{Objetivo General} desarrollar un generador
híbrido de historias con LLMs que combine un plan validado con la simulación de personajes con
memoria propia, junto con una metodología para evaluarlo basada en recuentos verificables y
comparaciones humanas por pares a ciegas. Para lograrlo se plantean los siguientes
\textbf{Objetivos Específicos}:

% Correspondencia objetivo -> capítulo: 1 -> Cap. 1 (Background.tex), 2 -> Cap. 2
% (Generator.tex), 3 -> Cap. 3 (Evaluation.tex), 4 -> Cap. 4 (Implementation.tex),
% 5 -> Cap. 5 (Experiments.tex). R08 escribe una conclusión por objetivo, en este orden.
\begin{itemize}
  \item \textbf{Analizar el estado del arte} de la generación de historias dirigida por la trama,
  por simulación de personajes e híbrida, y de su evaluación, para identificar qué adoptar y qué
  hueco queda. % -> Capítulo 1
  \item \textbf{Diseñar un generador híbrido} en el que un plan validado decida qué pasa, una
  función con personajes de memoria propia decida cómo pasa y un narrador decida cómo se
  cuenta. % -> Capítulo 2
  \item \textbf{Diseñar una metodología de evaluación} que combine rasgos contables con cita
  verificada, comparaciones humanas por pares a ciegas que incluyan una línea base de un solo
  \textit{prompt}, y un juez aprendido de esas comparaciones y validado fuera de
  muestra. % -> Capítulo 3
  \item \textbf{Implementar el generador y las herramientas de evaluación} en un sistema
  reproducible, con una interfaz para la lectura a ciegas. % -> Capítulo 4
  \item \textbf{Evaluar experimentalmente el generador}: validarlo antes de fijar su configuración,
  validar la extracción de rasgos, comparar parametrizaciones, recoger las preferencias de los
  lectores y ordenar el corpus de historias. % -> Capítulo 5
\end{itemize}

Las \textbf{hipótesis} de trabajo son tres. Las dos primeras comparan la configuración del
generador con una alternativa sobre la misma entrada; la tercera, cada mecanismo opcional con su
ausencia. Todas se contrastan con recuentos de rasgos y con la preferencia de los lectores en
comparaciones a ciegas, según los tres criterios de la evaluación: creatividad, desarrollo de los
personajes e interés de la trama.

% Correspondencia hipótesis -> experimentos (R06 y R07 en ../write_thesis_roadmap.md):
% H1 -> R06: celdas M03 (narrativo --plan-from) frente a su base M01 en la matriz A (EXP-1),
%       parcial: 14 de 35 unidades (P05 cerrada el 2026-10-30). R07: pares narrativo-simulado
%       del estudio humano (solo el piloto; P06 cerrada el 2026-10-31 sin estudio real).
% H2 -> R07: EXP-2 (--actor-memory own frente a shared) no se corrió (idea-bank desde P05.1,
%       2026-10-30); respuesta prevista: sin datos suficientes. Evidencia indirecta en R06: la
%       auditoría de 62 contextos de actor de la validación 7.5 (aislamiento, no mejora).
% H3 -> R06: validaciones antes de congelar el generador, P02.1 (narración, empate 2-2-1 en 5
%       pares) y P02.2 (brazos Top-Down, A 3 - B 1 en 4 pares), con un solo lector; R07, si el
%       juez o el piloto añaden algo.
% R12 (2026-11-03): las tres se responden en Experiments.tex, sec:respuesta-hipotesis, «sin datos
% suficientes» (en H3, lo observado es coherente con lo que anticipa: P02.1, P02.2 y 24-24 frente a la línea base en el piloto).
\begin{itemize}
  \item \textbf{H1:} con un mismo plan, la historia narrada a partir de una función difiere de la
  escrita directamente desde el plan en rasgos contables de personajes y de trama, más allá del
  ruido entre repeticiones, y los lectores la prefieren a ciegas en al menos uno de los tres
  criterios.
  \item \textbf{H2:} dar a cada personaje solo la memoria de lo que presenció, en lugar de la de
  todo lo ocurrido en público, produce escenas con menos menciones a hechos que el personaje no
  podía conocer, y que los lectores prefieren a ciegas.
  \item \textbf{H3:} más mecanismos no implican una historia mejor: que un mecanismo mejore medidas
  de la prosa no basta para que los lectores prefieran a ciegas las historias que lo usan frente a
  las que no. Traslada parte de la conclusión de WSE-bench~\cite{chen2026b}, que mide la coherencia
  con jueces LLM, a la preferencia de los lectores en este generador, y obliga a que cada mecanismo se
  gane su complejidad.
\end{itemize}

Las hipótesis no fijan de antemano qué resultado las confirmaría o refutaría. La regla de ruido de la
Sección~\ref{sec:horizontal} decide solo si una diferencia de la matriz se interpreta para un rasgo,
no cuántos rasgos tendrían que superarla para dar por cumplida H1, y la parte humana de las
hipótesis no tiene umbral.

El documento se organiza en cinco capítulos. El Capítulo~\ref{chapter:state-of-the-art} revisa el
estado del arte de la generación y la evaluación de historias y fija la posición del trabajo. El
Capítulo~\ref{chapter:generator} describe el generador híbrido: el plan, la función, el narrador y
los mecanismos opcionales. El Capítulo~\ref{chapter:evaluation} expone la metodología de
evaluación, y el Capítulo~\ref{chapter:implementation}, la implementación del sistema y sus
herramientas. El Capítulo~\ref{chapter:experiments} presenta los experimentos, con su alcance y sus
resultados incompletos, y responde a las hipótesis con lo que los datos permiten. Por último, las
Conclusiones recogen lo alcanzado en cada objetivo, las aportaciones y las limitaciones, y las
Recomendaciones, las líneas de trabajo futuro. Los anexos dan la lista núcleo de rasgos, las
peticiones de las obras de los experimentos y dos tablas complementarias.
\end{introduction}
